\section{A single separating Frobenius and a uniform trace bound for Case V}
\label{sec:case-v-host-number}

% BEGIN VUN1
\paragraph{VUN1. Restricting a host point to the Gaussian field.}
\label{claim:VUN1}
Let $K$ be any CM field satisfying $\mathrm{IsQi}$: it has degree
two over $\mathbb Q$ and contains a root of $X^2+1$. Work with the
source data $\Omega,\ell,c,N,e,\psi,o_1,S,(A_v),F,\iota,g_1$
in $\mathrm{HostNumber}_s$, including $\mathrm{HostData}$,
$\mathrm{IsHostFrame}$ and $F.\mathrm{Marked}(g_1)$. Write
$O=c.O$, $k=O/\mathfrak m_O$, $\Delta=G_{\mathbb Q}$,
$\pi:\Delta\twoheadrightarrow\Gamma=G_S$ and $H=\ker\eta$ in
$\Gamma$. Put $L=\overline{\mathbb Q}_\ell$ and write $\varepsilon$
for the cyclotomic character. The frame identities are
$\eta\pi=\eta_4$ and $\Psi\pi=\varepsilon^e\psi_{\rm Gal}$.

Take any marked characteristic-zero point $r$ of type I satisfying
the source $B'_s(K,\iota)$ conditions. Its automorphy clause gives
an embedding $j:K\hookrightarrow\Omega$, a representation
$\rho:G_K\to\mathrm{GL}_3(L)$ with matrices equal to those of
$r\pi$ on $G_K$, and the signed automorphy assertion for
$\rho\varepsilon^{-1}$ (HSR8). Both possible $j$ have image
$\mathbb Q(i)$ inside $\Omega$. Thus restriction identifies
$G_K$ with $\ker\eta_4\subset\Delta$, and its composite with
$\pi$ maps onto $H$. Indeed every lift to $\Delta$ of $h\in H$
has trivial $\eta_4$. A $\rho$-invariant subspace is invariant
under all the matrices on $H$, so type I gives irreducibility
over $L$. For fixed $j$, any isomorphism between two such
restrictions also intertwines all matrices on $H$.

The first conjunct of the supplied $\mathrm{IsAutCsdA}$ is
$\mathrm{AutCsd0}$. By HSR8 it consists exactly of a regular
algebraic cusp, its conjugate self-duality and the same attached
system as in $\mathrm{IsAutCsd}$, with twist $\varepsilon^{-1}$.
The matrices of $r$ are continuous, since a torus conjugate comes
from an adically continuous marked point over a finite coefficient
DVR (HSR1). Thus the matrices of $\rho$ and their inverses are
continuous, the latter also by $\rho(g)^{-1}=\rho(g^{-1})$.
For the $\mathrm{Adm3}$ branch of RNC16 we have proved the
irreducibility and automorphy conditions and the continuity field
of the level. We next prove its two local conditions at the
unchanged level $N$.
% END VUN1

% BEGIN VUN2
\paragraph{VUN2. The rank-three level on inertia away from $\ell$.}
\label{claim:VUN2}
Put $d_0=1,d_1=0,d_2=-1$, and on $G_K$ put
$\xi_i=(\psi_{\rm Gal}|_{G_K})^{d_i}$, embedded in $L^\times$.
These three global characters kill the group fixing $\mu_N$.
For every finite place $w\nmid\ell$ of $K$ there is an extension
of $w$ to $\Omega$ whose inertia under $\rho$ is simultaneously
diagonal in one basis, with these three characters.

Here are the local comparisons. Work first at the chosen $A_v$ over
a rational prime $p\ne\ell$ belonging to $S$, in an integral
model $r_D$ of the characteristic-zero point, with coefficient
DVR $D$. Its image over $L$ is a fixed torus conjugate of $r$.
On sign-positive inertia $I^+=I_v\cap\ker(\eta\pi)$ its
multiplier is one, and the base matrix is
$T(\Psi)=\operatorname{diag}(\Psi,1,\Psi^{-1})$.
At $v\in S_b$ the actual diagonal local basis of the integral
point gives exactly $T(\Psi)$ on $I^+$. At
$v\notin T_b=\{v_\ell\}\cup S_b$ both $r_D$ and $T(\Psi)$
kill $N_v=I^+\cap\ker\bar\Psi$ (CVI2). Their restrictions
therefore descend to matrices $U(t),V(t)$ on
$C=I^+/N_v\subseteq k^\times$, where $V=T(\Psi)$.
They agree upon reduction modulo $\mathfrak m_D$, and $|C|$ is
a unit in $D$. The matrix
\[
                h=|C|^{-1}\sum_{t\in C}U(t)V(t)^{-1}
\]
reduces to the identity, hence is invertible. Permuting the sum by
left multiplication gives $U(t)h=hV(t)$ for each $t$. This
constructs one basis valid on the entire group, including at the
dyadic prime. In either case $\Psi=\psi_{\rm Gal}$ on this
inertia because $\varepsilon$ is unramified at $p$.

For an arbitrary place $w$ above $p$, choose a prolongation $B$
to $\Omega$. If $p\notin S$, the image of its absolute inertia
in $\Gamma$ is trivial, so one may use the three trivial
characters there in the level condition; the displayed $\xi_i$
also restrict trivially since $p\nmid N$. Otherwise $B$ is
conjugate over $\mathbb Q$ to $A_v$. Its inertia over $K$
embeds in the conjugate of $I^+$, because $G_K$ is the sign
kernel. Transport the preceding basis, together with the torus
conjugation from the integral model. For a positive-sign
conjugating element this is ordinary matrix conjugation. For a
negative-sign element the extra involution on a matrix $X$ is
$X\mapsto JX^{-t}J$, with $J$ reversing the coordinates. It
fixes $T(\Psi)$, and $\Psi$ is invariant under conjugation as a
character. Thus the conjugated inertia representation still has
the indicated diagonal in one basis on this same $B$. This
establishes every field of the native $\mathrm{LevelAway}$.
% END VUN2

% BEGIN VUN3
\paragraph{VUN3. The ordinary flag and its cyclotomic inequalities.}
\label{claim:VUN3}
Above $\ell$ the required weight vector is $(e,0,-e)$. On
$D_{v_\ell}\cap\ker(\eta\pi)$, the integral local basis makes
$r_D$ upper triangular with diagonal $\Psi^{d_i}$ on inertia.
Choose any $w\mid\ell$ of $K$ and any prolongation $B$ to $\Omega$.
Transport a local basis from the stored rational prolongation to $B$,
using conjugacy over $\mathbb Q$ as in VUN2. The group
$D_{K,B}$ identifies, inside $\Delta$, with
$D_{\mathbb Q,B}\cap\ker\eta_4$. Under a negative conjugation
$JX^{-t}J$ again preserves upper triangularity; it reverses and
inverts the three diagonal characters. On inertia this still
gives $T(\Psi)$, since this torus is fixed. Thus $\rho$ on
$D_{K,B}$ has one invariant complete ordered flag with characters
$\kappa_i$. On its inertia they satisfy
\[
                \kappa_i=\varepsilon^{e d_i}\xi_i
\]
with the same three global level-$N$ characters as in VUN2.

For the last disjunction in $\mathrm{LevelAt}$, we prove that
$\kappa_i\ne\kappa_j\varepsilon$ for every $i,j$, including
$e=1$. HSR3 gives a sign-positive element in the stored
decomposition group with $\varepsilon(g)=1$ and $\Psi(\pi g)^2-1$
a unit. Its conjugate at $B$ lies in $D_{K,B}$. At this element
the diagonal entries on the flag are units with pairwise distinct
residues: HSR4 proves this from the three distinct residual roots
of the integral characteristic polynomial. Replacing the three
roots by their inverses in reverse order preserves this property.
For $i\ne j$ they therefore differ at an element where
$\varepsilon=1$. For $i=j$, equality would force $\varepsilon$
trivial on $D_{K,B}$. This is impossible since the finite local
extension $K_w/\mathbb Q_\ell$ cannot contain all $\ell$-power
roots of unity, whose local degrees grow without bound.
This proves the unlinked alternative, independently of any
potential-crystallinity alternative.

Together VUN1--3 show that every $\rho$ supplied by the automorphy
witness of an eligible point satisfies the literal
$\mathrm{Adm3}(\iota,e,N,\rho)$ of RNC16. This specialization
uses the native point's local inertial basis and the host separator;
its local conclusions apply to both embeddings of the Gaussian field.
% END VUN3

% BEGIN VUN4
\paragraph{VUN4. One prime and one Frobenius element for every point.}
\label{claim:VUN4}
Take $m=\ell N$. Some prime divisor $p$ of $12m-1$ is $3$ modulo
$4$: the integer is $3$ modulo $4$, whereas a product of primes
$1$ modulo $4$ cannot have this residue. The prime $p$ divides
neither $3$ nor $m$. Hence $7\le p\le 12\ell N$ and $p\nmid
\ell N$. This prime is inert in $K\simeq\mathbb Q(i)$:
$X^2+1$ has no root in $\mathbb F_p$ but splits in $\mathbb F_{p^2}$.
The unique prime $w\mid p$ of $K$ has norm $q=p^2$, and the
discriminant of $K$ is $-4$, so $p\nmid\operatorname{disc} K$.

Fix a prolongation $B$ of the $p$-adic valuation to $\Omega$.
In the decomposition group over $\mathbb Q$ choose a lift of
arithmetic Frobenius and square it to obtain $\sigma$. Then
$\varepsilon(\sigma)=p^2$ and $\eta_4(\sigma)=1$.
For \emph{each} embedding $j:K\hookrightarrow\Omega$ the same
$\sigma$ fixes its image and defines $\tau_j\in G_K$. It belongs
to the decomposition group at $B$ and acts on residue elements
as the $p^2$-power map; it is an arithmetic Frobenius at $w$.
The two embeddings induce the same $w$ since $p$ is inert. Put
$\phi=\pi(\sigma^{-1})$. This one element of $\Gamma$ is fixed
before any representation or finite collection, even when
automorphy witnesses use different embeddings.

Set $z_0=\iota(\psi_{\rm Gal}(\sigma^{-1}))$, with
$O\hookrightarrow L$ understood, and $d=\varphi(N)$. The
Dirichlet character factors through $(\mathbb Z/N)^\times$, so
$z_0^d=1$. At the inverse Frobenius
$\Psi(\phi)=q^{-e}\psi_{\rm Gal}(\sigma^{-1})$. The three
diagonal terms of the original base point now give
\[
 \operatorname{tr}(r_z(\phi)^{\rm mat})
   = \iota^{-1}\!\left(q^{-e}z_0+1+q^e z_0^{-1}\right).
\]
This is precisely the geometric Frobenius convention of RNC16.
% END VUN4

% BEGIN VUN5
\paragraph{VUN5. The integer bound for one embedding of $K$.}
\label{claim:VUN5}
Fix the constants $a,b\in\mathbb N$ furnished by RNC16 for $K$,
independently of $\ell,c,N,e,\psi$ and the frame. At the prime
of VUN4 use the orbit $\mathcal O=\mu_d(\mathbb C)$, with each
root listed once and $d=\varphi(N)\ge1$. Every complex field
automorphism permutes this indexed orbit. The other catalogue
conditions hold: $p\ge7$, $p\nmid N\operatorname{disc}K$,
$\ell\ne p$ and $\ell$ is odd by $\mathrm{HostData}$.
Consequently RNC16 gives an integer $M\ne0$ with
\[
             \log |M|\le a\varphi(N)[N(e+1)]^b\log p .
\]
Its nonvanishing and product estimate hold at every subsequent
allowed coefficient realization. Here $p$ was selected using
$\ell$; this is an allowed specialization, since for this $p$
the catalogue works at every odd prime different from $p$.

Fix one embedding $j$ and a finite indexed family of eligible
points with automorphy witnessed through $j$. Assume that an
isomorphism of two members' matrices on $H$ implies equality
of their indices. Use the supplied restrictions $\rho_i$,
each satisfying the rank-three branch of the catalogue by VUN3.
Put their images into the one set $X_{z_0}$ and set all other
$X_z$ empty. By VUN1 an isomorphism of two restrictions would
identify the indices, so passing from the indexed family to its
image loses no factors. For each $i$ the value of $\rho_i$ at
$\tau_j^{-1}$ is exactly the matrix of $r_i$ at $\phi$.
By VUN4 the comparison trace in the catalogue is the original
base trace. Write $\delta_i=\operatorname{tr}(r_i(\phi)^{\rm mat})-
\operatorname{tr}(r_z(\phi)^{\rm mat})$. RNC16 gives all
$\delta_i\ne0$ and
\[
                |M|_\ell\le\prod_i|\delta_i|_\ell .
\]
The singleton case requires no prior pairwise hypothesis. In
particular the same $\phi$ separates \emph{every} allowed marked
type-I point satisfying $B'_s$.

Normalize $v_\ell(\ell)=1$. Since $M\ne0$ is an integer,
$\ell^{v_\ell(M)}$ divides $|M|$. Applying the negative
logarithm to the product bound therefore gives
\[
 \log\ell\sum_i v_\ell(\delta_i)
    \le -\log|M|_\ell
    \le \log |M|
    \le a\varphi(N)[N(e+1)]^b\log p.
\]
The product and sum conventions also give the stated inequality
for an empty fiber.
% END VUN5

% BEGIN VUN6
\paragraph{VUN6. The literal uniform host number.}
\label{claim:VUN6}
For the Gaussian field $K$, choose the exponent
$C=2a+b+4$ from VUN5 before all closures, primes, coefficient
rings, levels, characters, frames, complex embeddings and marks.
For any such data take $\phi$ from VUN4. For every later finite
list of marked characteristic-zero type-I points satisfying
$B'_s(K,\iota)$ and the precise pairwise nonisomorphism test
on $H$, choose one automorphy witness for each. Partition by the
associated field embedding $K\hookrightarrow\Omega$. There are
two embeddings, so summing VUN5 over at most two fibers gives
\[
             \log\ell\sum_i v_\ell(\delta_i)
               \le 2a\varphi(N)[N(e+1)]^b\log p .
\]
Each application uses just the one embedding of $K$ attached
to that fiber.

Put $X=\ell N e$. The host has $\ell\ge3$, $N$ positive and
even, and $e\ge1$, hence $X\ge6$. We have $2a\le X^{2a}$
(also for $a=0$), $\varphi(N)\le X$, $N(e+1)\le X$, and
$\log p\le p\le12\ell N\le X^3$. Multiplication proves
\[
  \left(\sum_i v_\ell\bigl(
     \operatorname{tr}(r_i(\phi)^{\rm mat})
      -\operatorname{tr}(r_z(\phi)^{\rm mat})\bigr)\right)\log\ell
       \le (\ell N e)^C .
\]
For all these nonzero arguments the source definition
$\mathrm{vl1}_\ell(x)=-\log|x|_\ell/\log\ell$ is precisely
$v_\ell(x)$. This inequality and the separate singleton trace
separation of VUN5 supply both clauses of the original
$\mathrm{HostNumber}_s$ for every marked $g_1$. The inherited
catalogue and automorphy comparisons are RNC16 and HSR8, used
on the native $\mathrm{Adm3}$ branch verified in VUN1--3.
% END VUN6
